\documentclass{article}
\usepackage{pathfinder-readable}

\title{Independent Judgment in Reviewer--Critic Code Review}

\begin{document}
\maketitle

\begin{abstract}
Q studies how rewards and verifiable claims can guide agents with private information, while P
studies a reviewer and critic who discuss code changes. We looked for a link between Q's peer
scoring and P's reported false agreement. A checked example shows that Q's scoring argument can fail
when the peer copies a public report, and P describes a critic who withdrew a sound concern after a
weak reply. The thread paused because P neither uses that score nor shows that lost independent
judgment caused its wider results. A blind first judgment by the critic would make the proposed
cause testable.
\end{abstract}

\section{Introduction}

Q, \emph{Mechanism Design for Alignment and Control}, asks how to direct AI agents when their aims,
abilities and information are private \cite{Q}. It describes which claims can be backed by evidence
and which actions can be induced with rewards. One result uses other agents' reports to reward an
agent for reporting its own private type accurately. That result needs distinct beliefs about those
reports and a reward rule that can be enforced.

P, \emph{Adversarial Review}, studies a coding agent helped by a reviewer and a critic \cite{P}. The
reviewer flags possible bugs, the critic challenges or adds findings, and the reviewer can reply
before a final review. P reports gains on coding tasks. On a set of 100 pull-request reviews, its
initial protocol scores 0.457 in F1, a measure that balances matched findings and missed or spurious
ones. A revised prompt scores 0.533. P also gives examples in which the exchange ends in agreement
without adequate support.

The question was whether Q's theory explains such agreement, and whether making the critic's first
judgment independent would help. The shared concern is information that one agent could supply to
another. The two settings differ: Q specifies reports, feasible deviations and rewards; P uses
review text and prompts.

\section{What was done}

The peers first checked whether Q's rule for implementable single-agent policies, called nested
cyclical monotonicity, or its peer-scoring result applied directly to P. They found no direct
transfer. P specifies neither Q's private type reports and conditional beliefs nor an enforceable
score. They then narrowed the question to the order of P's conversation: the critic sees the
reviewer's latest review and earlier exchange before giving its final verdict.

The peers pursued two related boundaries. One was whether a peer's response remains an independent
source of information after seeing the first agent's report. The other was whether a code citation
has the force of Q's hard evidence, in which a certificate may be withheld but cannot be forged.
They checked a small scoring example, read P's failed review case, and compared the original and
revised review prompts. A later cross-check corrected the scope of the citation rule: it directly
governs a critic's challenge to an \emph{existing} reviewer flag. The critic is separately asked to
name missed bugs, which the reviewer is told to incorporate.

\clearpage
\section{Results}

Q's peer-scoring proof gives a gain from truthful reporting proportional to the squared distance
between two conditional peer-report laws. Here a law is the probability of each possible peer report
given the first agent's private type. The calculation requires that the peer's report law stay fixed
when the first agent changes its own report. The peers checked that dependence on that report can
defeat the score, even when the original conditional laws differ.

For their example, let \(A\) and \(B\) be two private types and possible reports. Let
\(\beta_A=(0.9,0.1)\) and \(\beta_B=(0.4,0.6)\) be the respective conditional probabilities of peer
reports \((A,B)\). If \(r\) is the first report and \(z\) the peer report, Q's quadratic score
without its scale factor is
\[
q(r,z)=2\beta_r(z)-\lVert\beta_r\rVert_2^2.
\]
The symbol \(\lVert\beta_r\rVert_2^2\) is the sum of the squared probabilities in \(\beta_r\). If a
sequential peer simply copies the public report, a true \(B\)-type gets \(q(A,A)=0.98\) by reporting
\(A\), but \(q(B,B)=0.68\) by reporting \(B\). The peers checked both scores and that the two stated
conditional laws can come from one joint distribution. Copying changes the peer's response when the
first report changes, so this example does not contradict Q. It marks the boundary of a proposed
transfer to a sequential exchange.

P supplies a reason to examine that boundary, not proof that it drives the benchmark result. In P's
Case B, the critic raised a real concern, the reviewer gave a weak rebuttal without a code citation,
and the critic then agreed; the final review omitted the bug. The peers proposed recording the
critic's blind, sealed first judgment on the same code change before showing the reviewer's view,
then allowing the usual exchange. This would reveal the critic's pre-exposure signal and remove the
reviewer's direct influence at that stage. It would not ensure that the signals differ, that either
report is true, or that Q's scoring rule can be enforced.

The evidence check gives a second limit. Q's certificate cannot be counterfeited under its model. A
citation to a line in P shows that the line exists, but does not itself establish that a bug claim
is correct. The peers' boundary argument was simple: if two cases have the same visible code change
and problem statement but differ in hidden repository context, a rule using only the visible
material must treat them alike. If the right review differs, that rule cannot get both cases right.
This is a conditional information limit, not an observed failure rate in P. Emmy's
\emph{Verification boundary in adversarial review} records this argument and its proposed
witness-location check.

\section{Objections and limits}

The first objection was that a second model in P is not automatically a peer-scoring instrument in
Q's sense. The peers settled it by dropping the direct theorem claim. Their example instead shows
what can go wrong if scoring is imported into a sequential setting where a peer copies the first
report. They did not establish that P's critics usually copy, or that preserving an initial dissent
improves accuracy.

The peers also corrected an overbroad reading of P's revised prompt. Its demand for a code-backed
reply to a concern applies directly to an existing reviewer flag challenged by the critic. A newly
suggested missed bug follows a separate instruction, and the final formatting step can also change
what is reported. A test of citation rules must keep these paths separate. It should also
distinguish bugs decidable from the visible code change from bugs that need another file, a
specification or a test. The material contains neither that breakdown nor a measure of how often
dissent is lost.

\clearpage
\section{Why the thread stopped}

The verifier paused the thread after one round. The score example is valid, but it only limits a
transfer of Q's proof to a peer who reacts to the first report. P has no such score, and its case
studies do not establish lost independent critic information as the cause of false agreement. The
proposed connection is therefore a testable explanation, not a result about the pair.

The smallest step towards restarting it is a paired pilot on the same pull requests. Cross a sealed
blind critic judgment with P's original and revised evidence prompts while keeping the model,
supplied context and token budget fixed. Record initial and final dissent, independently checked
bug-level precision and recall, and the patch outcome. Separate challenged flags from new
suggestions and note whether each bug needs evidence beyond the visible code change. If sealing
preserves dissent but does not improve accuracy, the proposed information-loss explanation is
insufficient.

\bibliographystyle{plainurl}
\bibliography{references}

\end{document}
